\chapter{مبانی نظری روش نیوتن}\label{ch:theory}

\section{مقدمه}

روش نیوتن یکی از بنیادی‌ترین و پرکاربردترین روش‌های عددی برای حل معادلات غیرخطی است. این روش بر پایه ترکیب دو ایده مهم بنا شده است: نخست، تقریب محلی یک تابع غیرخطی به وسیله خط مماس و دوم، تولید دنباله‌ای از تقریب‌ها که به ریشه تابع نزدیک شوند. به همین دلیل، روش نیوتن را می‌توان پلی میان مفاهیم حساب دیفرانسیل و محاسبات عددی دانست.

در فصل قبل بیان شد که بسیاری از مسائل علمی و مهندسی در نهایت به حل معادله‌ای از فرم زیر منجر می‌شوند:
\begin{equation}
f(x)=0
\label{eq:ch2-fzero}
\end{equation}
در این رابطه، هدف یافتن مقداری از $x$ است که مقدار تابع را صفر کند. اگر چنین مقداری را با $\alpha$ نمایش دهیم، آنگاه داریم:
\begin{equation}
f(\alpha)=0
\label{eq:ch2-falpha}
\end{equation}

در بسیاری از موارد، محاسبه دقیق $\alpha$ امکان‌پذیر نیست. بنابراین به جای جواب دقیق، به دنبال یک تقریب مناسب از آن هستیم. روش نیوتن با شروع از یک حدس اولیه، دنباله‌ای از مقادیر تولید می‌کند که در شرایط مناسب به ریشه واقعی همگرا می‌شود.

در این فصل، ابتدا مفهوم ریشه تابع و مسئله ریشه‌یابی توضیح داده می‌شود. سپس مفهوم تقریب عددی و خطا بررسی می‌گردد. در ادامه، روش‌های تکراری و معیارهای توقف معرفی می‌شوند. پس از آن، فرمول روش نیوتن از دو دیدگاه هندسی و تحلیلی استخراج می‌شود. در نهایت، شرایط همگرایی و محدودیت‌های روش مورد بررسی قرار می‌گیرد.

\section{ریشه تابع و مسئله ریشه‌یابی}
\label{sec:root}

یکی از مسائل مهم در ریاضیات کاربردی، یافتن نقاطی است که در آن‌ها مقدار یک تابع برابر صفر می‌شود. اگر $f$ تابعی حقیقی باشد، عدد $\alpha$ را ریشه تابع $f$ می‌نامیم هرگاه رابطه \eqref{eq:ch2-falpha} برقرار باشد.

به بیان هندسی، ریشه‌های تابع نقاطی هستند که نمودار تابع محور افقی را قطع می‌کند یا بر آن مماس می‌شود. برای مثال، اگر تابع $f(x)=x^2-4$ را در نظر بگیریم، ریشه‌های آن برابرند با:
\[
x=-2
\]
و
\[
x=2
\]
زیرا:
\[
f(-2)=(-2)^2-4=0
\]
و
\[
f(2)=2^2-4=0
\]

در حالت کلی، یک تابع ممکن است هیچ ریشه حقیقی نداشته باشد، یک ریشه داشته باشد یا دارای چندین ریشه باشد. به عنوان نمونه، تابع $f(x)=x^2+1$ در مجموعه اعداد حقیقی ریشه‌ای ندارد؛ زیرا مقدار آن برای هر $x$ حقیقی مثبت است. در مقابل، تابع $f(x)=\sin x$ دارای بی‌نهایت ریشه است.

مسئله ریشه‌یابی معمولاً به این صورت بیان می‌شود: تابعی مانند $f$ داده شده است و باید مقداری مانند $\alpha$ پیدا شود به‌طوری‌که \eqref{eq:ch2-falpha} برقرار باشد.

در روش‌های عددی، معمولاً به جای یافتن مقدار دقیق $\alpha$، به دنبال عددی مانند $x_n$ هستیم که مقدار آن به اندازه کافی به $\alpha$ نزدیک باشد. به عبارت دیگر، می‌خواهیم:
\[
x_n \approx \alpha
\]
و در نتیجه مقدار تابع در $x_n$ نیز تا حد قابل قبولی به صفر نزدیک شود:
\[
f(x_n)\approx 0
\]
این دیدگاه، پایه اصلی روش‌های عددی ریشه‌یابی است.

\section{تقریب عددی و مفهوم خطا}

در محاسبات عددی، معمولاً جواب دقیق مسئله در دسترس نیست یا محاسبه آن دشوار است. بنابراین به جای مقدار دقیق، از یک مقدار تقریبی استفاده می‌شود. برای ارزیابی کیفیت این مقدار تقریبی، مفهوم خطا تعریف می‌شود.

فرض کنید $\alpha$ ریشه واقعی معادله \eqref{eq:ch2-fzero} و $x_n$ تقریب عددی آن در مرحله $n$-ام باشد. خطای واقعی در این مرحله به صورت زیر تعریف می‌شود:
\begin{equation}
e_n = |\alpha - x_n|
\label{eq:ch2-abs-error}
\end{equation}

این رابطه نشان می‌دهد که هرچه مقدار $e_n$ کوچک‌تر باشد، تقریب $x_n$ به ریشه واقعی نزدیک‌تر است. با این حال، در بسیاری از مسائل عملی، مقدار دقیق $\alpha$ معلوم نیست؛ بنابراین خطای واقعی نیز مستقیماً قابل محاسبه نیست.

به همین دلیل، معمولاً از معیارهایی استفاده می‌شود که بدون دانستن ریشه واقعی قابل محاسبه باشند. یکی از رایج‌ترین معیارها، اختلاف دو تقریب متوالی است:
\begin{equation}
|x_{n+1}-x_n|
\label{eq:ch2-diff}
\end{equation}

اگر این مقدار کوچک باشد، می‌توان نتیجه گرفت که دنباله تقریب‌ها به پایداری نسبی رسیده است. معیار دیگر، مقدار تابع در تقریب فعلی است:
\begin{equation}
|f(x_n)|
\label{eq:ch2-fval}
\end{equation}

اگر مقدار $|f(x_n)|$ نزدیک صفر باشد، آنگاه $x_n$ یک تقریب مناسب برای ریشه محسوب می‌شود.

در بسیاری از الگوریتم‌های عددی، مقدار کوچکی مانند $\varepsilon$ به عنوان آستانه خطا انتخاب می‌شود. الگوریتم زمانی متوقف می‌شود که یکی از شرایط زیر برقرار باشد:
\begin{equation}
|x_{n+1}-x_n|<\varepsilon
\label{eq:ch2-stop-diff}
\end{equation}
یا:
\begin{equation}
|f(x_n)|<\varepsilon
\label{eq:ch2-stop-func}
\end{equation}

انتخاب مقدار $\varepsilon$ به دقت مورد نیاز مسئله بستگی دارد. برای مثال، اگر دقت بالاتری لازم باشد، مقدار $\varepsilon$ کوچک‌تر انتخاب می‌شود. البته انتخاب مقدار بسیار کوچک برای $\varepsilon$ ممکن است باعث افزایش تعداد تکرارها و هزینه محاسباتی شود.

گاهی برای بررسی دقیق‌تر، از خطای نسبی نیز استفاده می‌شود. خطای نسبی میان دو تقریب متوالی را می‌توان به صورت زیر تعریف کرد:
\begin{equation}
\frac{|x_{n+1}-x_n|}{|x_{n+1}|}
\label{eq:ch2-rel-error}
\end{equation}

این معیار زمانی مفید است که اندازه مقدار جواب در ارزیابی خطا اهمیت داشته باشد. برای جلوگیری از تقسیم بر صفر، در برخی پیاده‌سازی‌ها از صورت اصلاح‌شده زیر استفاده می‌شود:
\begin{equation}
\frac{|x_{n+1}-x_n|}{\max(1,|x_{n+1}|)}
\label{eq:ch2-rel-error-mod}
\end{equation}

در نتیجه، مفهوم خطا نقش اساسی در روش‌های عددی دارد؛ زیرا هم کیفیت جواب تقریبی را مشخص می‌کند و هم معیار توقف الگوریتم را فراهم می‌سازد.

\section{روش‌های تکراری و معیار توقف}

روش نیوتن از دسته روش‌های تکراری است. در یک روش تکراری، ابتدا یک مقدار اولیه برای جواب انتخاب می‌شود و سپس با استفاده از یک رابطه مشخص، مقادیر جدیدی تولید می‌گردد. این فرایند تا زمانی ادامه می‌یابد که جواب به دقت مورد نظر برسد.

به طور کلی، یک روش تکراری را می‌توان به شکل زیر نمایش داد:
\begin{equation}
x_{n+1}=g(x_n)
\label{eq:ch2-iteration}
\end{equation}

در این رابطه، $x_n$ تقریب فعلی و $x_{n+1}$ تقریب بعدی است. تابع $g$ را تابع تکرار می‌نامند. اگر دنباله $\{x_n\}$ به عددی مانند $\alpha$ نزدیک شود، گفته می‌شود روش به $\alpha$ همگرا شده است:
\begin{equation}
\lim_{n\to\infty} x_n = \alpha
\label{eq:ch2-convergence}
\end{equation}

در چنین حالتی، $\alpha$ نقطه ثابت تابع $g$ خواهد بود؛ یعنی:
\begin{equation}
\alpha = g(\alpha)
\label{eq:ch2-fixedpoint}
\end{equation}

در روش‌های عددی، از آنجا که نمی‌توان بی‌نهایت تکرار انجام داد، باید معیار مشخصی برای توقف الگوریتم در نظر گرفت. معیار توقف نشان می‌دهد چه زمانی تقریب فعلی به اندازه کافی دقیق است.

برای روش نیوتن، معیارهای توقف رایج عبارت‌اند از:
\begin{equation}
|x_{n+1}-x_n|<\varepsilon
\label{eq:ch2-stop-diff2}
\end{equation}
\begin{equation}
|f(x_{n+1})|<\varepsilon
\label{eq:ch2-stop-func2}
\end{equation}
و:
\begin{equation}
n \geq N_{\max}
\label{eq:ch2-stop-maxiter}
\end{equation}

که در آن، $N_{\max}$ بیشینه تعداد تکرارهای مجاز است. شرط سوم برای جلوگیری از اجرای بی‌پایان الگوریتم به کار می‌رود. زیرا در برخی موارد ممکن است روش نیوتن همگرا نشود یا بسیار کند همگرا شود.

بنابراین، یک معیار توقف کامل می‌تواند ترکیبی از چند شرط باشد. برای مثال، الگوریتم زمانی متوقف شود که یکی از شرایط \eqref{eq:ch2-stop-diff2} یا \eqref{eq:ch2-stop-func2} برقرار گردد یا تعداد تکرارها از حد مجاز بیشتر شود.

این معیارها در پیاده‌سازی عملی روش نیوتن اهمیت زیادی دارند، زیرا هم دقت جواب و هم هزینه محاسباتی را کنترل می‌کنند.

\section{ایده هندسی روش نیوتن}

روش نیوتن از نظر هندسی بسیار قابل فهم است. فرض کنید تابع $f$ دارای ریشه‌ای در نزدیکی نقطه $x_n$ باشد. اگر مقدار $x_n$ به عنوان تقریب فعلی ریشه در نظر گرفته شود، می‌توان در نقطه $(x_n,f(x_n))$ خط مماس بر نمودار تابع را رسم کرد.

از آنجا که خط مماس، رفتار تابع را در نزدیکی نقطه $x_n$ تقریب می‌زند، می‌توان انتظار داشت که محل برخورد این خط با محور $x$، تقریب بهتری برای ریشه تابع باشد. این محل برخورد را $x_{n+1}$ می‌نامیم.

به بیان دیگر، در روش نیوتن به جای حل مستقیم معادله غیرخطی \eqref{eq:ch2-fzero}، در هر مرحله، یک معادله خطی تقریبی حل می‌شود. این معادله خطی همان معادله خط مماس بر نمودار تابع در نقطه فعلی است.

معادله خط مماس در نقطه $x_n$ به صورت زیر است:
\begin{equation}
y=f(x_n)+f'(x_n)(x-x_n)
\label{eq:ch2-tangent}
\end{equation}

حال برای یافتن محل برخورد این خط با محور افقی، مقدار $y=0$ قرار داده می‌شود:
\begin{equation}
0=f(x_n)+f'(x_n)(x-x_n)
\label{eq:ch2-tangent-zero}
\end{equation}

مقدار به‌دست‌آمده برای $x$ همان تقریب بعدی $x_{n+1}$ است.

از دیدگاه تصویری، این فرایند به شکل زیر قابل توصیف است:
\begin{enumerate}
    \item یک حدس اولیه $x_0$ انتخاب می‌شود.
    \item خط مماس بر نمودار تابع در نقطه $x_0$ رسم می‌شود.
    \item محل برخورد مماس با محور $x$، نقطه $x_1$ را مشخص می‌کند.
    \item همین فرایند برای $x_1$، سپس $x_2$ و مراحل بعدی تکرار می‌شود.
\end{enumerate}

اگر تابع رفتار مناسبی داشته باشد و حدس اولیه به ریشه نزدیک باشد، این دنباله معمولاً با سرعت زیادی به ریشه واقعی نزدیک می‌شود.

\section{استخراج تحلیلی فرمول روش نیوتن}
\label{sec:formula}

برای استخراج فرمول روش نیوتن، از معادله خط مماس \eqref{eq:ch2-tangent} استفاده می‌کنیم. فرض کنید $x_n$ تقریب فعلی ریشه معادله \eqref{eq:ch2-fzero} باشد. در روش نیوتن، تقریب بعدی ریشه، یعنی $x_{n+1}$، نقطه‌ای است که این خط مماس محور افقی را قطع می‌کند. بنابراین مقدار $y$ را صفر قرار می‌دهیم:
\begin{equation}
0=f(x_n)+f'(x_n)(x_{n+1}-x_n)
\label{eq:ch2-newton-deriv1}
\end{equation}

حال این رابطه را نسبت به $x_{n+1}$ حل می‌کنیم:
\[
f'(x_n)(x_{n+1}-x_n)=-f(x_n)
\]

پس:
\begin{equation}
x_{n+1}-x_n=-\frac{f(x_n)}{f'(x_n)}
\label{eq:ch2-newton-step}
\end{equation}

و در نتیجه فرمول اصلی روش نیوتن به دست می‌آید:
\begin{equation}
x_{n+1}=x_n-\frac{f(x_n)}{f'(x_n)}
\label{eq:ch2-newton}
\end{equation}

این فرمول که در بسیاری از کتاب‌های استاندارد آنالیز عددی مانند \cite{burden2010numerical} به تفصیل بررسی شده است، زمانی قابل استفاده است که مشتق تابع در نقطه $x_n$ صفر نباشد؛ یعنی $f'(x_n)\neq 0$. زیرا اگر $f'(x_n)=0$ باشد، کسر موجود در \eqref{eq:ch2-newton} تعریف‌نشده می‌شود.

برای نمایش دقیق‌تر این موضوع می‌توان رابطه نیوتن را به صورت چندضابطه‌ای نوشت:
\begin{equation}
x_{n+1}=
\begin{cases}
x_n-\dfrac{f(x_n)}{f'(x_n)} & \text{اگر } f'(x_n)\neq 0,\\[8pt]
\text{تعریف‌نشده} & \text{اگر } f'(x_n)=0.
\end{cases}
\label{eq:ch2-newton-piecewise}
\end{equation}

این رابطه نشان می‌دهد که روش نیوتن تنها در حالتی قابل ادامه است که مشتق تابع در نقطه فعلی صفر نباشد.

\section{تفسیر روش نیوتن بر اساس بسط تیلور}

روش نیوتن را می‌توان از دیدگاه بسط تیلور نیز تحلیل کرد. این تفسیر نشان می‌دهد که روش نیوتن در واقع از تقریب خطی تابع در نزدیکی نقطه فعلی استفاده می‌کند.

اگر تابع $f$ در اطراف نقطه $x_n$ مشتق‌پذیر باشد، بسط تیلور مرتبه اول تابع در نزدیکی $x_n$ به صورت زیر نوشته می‌شود:
\begin{equation}
f(x)\approx f(x_n)+f'(x_n)(x-x_n)
\label{eq:ch2-taylor1}
\end{equation}

هدف ما یافتن مقداری از $x$ است که $f(x)=0$. بنابراین در تقریب خطی بالا، مقدار سمت چپ را صفر در نظر می‌گیریم:
\begin{equation}
0\approx f(x_n)+f'(x_n)(x-x_n)
\label{eq:ch2-taylor1-zero}
\end{equation}

با حل این رابطه برای $x$، فرمول نیوتن \eqref{eq:ch2-newton} به دست می‌آید. بنابراین، روش نیوتن در هر مرحله تابع غیرخطی را با یک تقریب خطی جایگزین کرده و ریشه این تقریب خطی را به عنوان تقریب جدید ریشه تابع اصلی انتخاب می‌کند.

برای تحلیل دقیق‌تر خطای روش، می‌توان از بسط تیلور مرتبه دوم استفاده کرد. فرض کنید $\alpha$ ریشه واقعی تابع باشد، یعنی $f(\alpha)=0$. اگر تابع $f$ دوبار مشتق‌پذیر باشد، بسط تیلور تابع حول $x_n$ را می‌توان به صورت زیر نوشت:
\begin{align}
f(\alpha) &=
f(x_n)+f'(x_n)(\alpha-x_n)
+\frac{f''(\xi_n)}{2}(\alpha-x_n)^2,
\label{eq:ch2-taylor2}\\
0 &=
f(x_n)+f'(x_n)(\alpha-x_n)
+\frac{f''(\xi_n)}{2}(\alpha-x_n)^2,
\label{eq:ch2-taylor2-zero}\\
\alpha - x_{n+1}
&=
-\frac{f''(\xi_n)}{2f'(x_n)}(\alpha-x_n)^2.
\label{eq:ch2-error-relation}
\end{align}
که در آن $\xi_n$ عددی بین $x_n$ و $\alpha$ است. رابطه \eqref{eq:ch2-error-relation} پایه‌ای برای تحلیل همگرایی روش نیوتن است و نشان می‌دهد که جمله خطای مرتبه دوم در رفتار روش نقش مهمی دارد.

\section{همگرایی روش نیوتن}
\label{sec:conv}

یکی از مهم‌ترین ویژگی‌های روش نیوتن، سرعت بالای همگرایی آن است. اگر تابع شرایط مناسبی داشته باشد و حدس اولیه به اندازه کافی به ریشه واقعی نزدیک باشد، روش نیوتن معمولاً با سرعت بسیار زیادی به جواب همگرا می‌شود.

فرض کنید $\alpha$ ریشه ساده تابع $f$ باشد؛ یعنی $f(\alpha)=0$ و $f'(\alpha)\neq 0$. همچنین فرض کنید تابع $f$ در نزدیکی $\alpha$ دوبار مشتق‌پذیر باشد. در این صورت، اگر حدس اولیه $x_0$ به اندازه کافی به $\alpha$ نزدیک باشد، دنباله تولیدشده توسط روش نیوتن به $\alpha$ همگرا می‌شود.

خطای مرحله $n$-ام را به صورت زیر تعریف می‌کنیم:
\begin{equation}
e_n=x_n-\alpha
\label{eq:ch2-error-def}
\end{equation}

در شرایط مناسب، رابطه خطا در روش نیوتن تقریباً به صورت زیر است و از مرتبه دوم همگرایی برخوردار است \cite{atkinson1989introduction}:
\begin{equation}
e_{n+1}\approx C e_n^2
\label{eq:ch2-quadratic-error}
\end{equation}
که در آن $C$ ثابتی وابسته به تابع و ریشه است. این رابطه نشان می‌دهد که خطای مرحله بعد تقریباً متناسب با مربع خطای مرحله فعلی است. این نوع همگرایی را همگرایی درجه دوم می‌نامند.

همگرایی درجه دوم بسیار سریع است. برای مثال، اگر خطای یک مرحله برابر $10^{-2}$ باشد، در مرحله بعد ممکن است در حدود $10^{-4}$ شود و در مرحله بعد به حدود $10^{-8}$ برسد. البته این رفتار زمانی رخ می‌دهد که تقریب‌ها وارد ناحیه همگرایی شده باشند.

به صورت دقیق‌تر، ثابت همگرایی برای روش نیوتن در نزدیکی ریشه ساده معمولاً به شکل زیر ظاهر می‌شود:
\begin{equation}
C=\left|\frac{f''(\alpha)}{2f'(\alpha)}\right|
\label{eq:ch2-conv-constant}
\end{equation}

در نتیجه می‌توان نوشت:
\begin{equation}
|e_{n+1}| \approx 
\left|\frac{f''(\alpha)}{2f'(\alpha)}\right|
|e_n|^2
\label{eq:ch2-error-approx}
\end{equation}

این رابطه نشان می‌دهد که اگر مشتق اول تابع در ریشه بسیار کوچک باشد یا مشتق دوم بزرگ باشد، رفتار همگرایی ممکن است تحت تأثیر قرار گیرد. رابطه \eqref{eq:ch2-error-relation} نیز همین مطلب را از دیدگاه بسط تیلور تأیید می‌کند.

\section{محدودیت‌ها و حالت‌های شکست روش نیوتن}

با وجود قدرت و سرعت بالای روش نیوتن، این روش همیشه موفق نیست. در برخی شرایط، روش ممکن است به ریشه همگرا نشود، به ریشه‌ای غیرمنتظره همگرا شود یا حتی دچار نوسان و واگرایی گردد. شناخت این محدودیت‌ها برای استفاده درست از روش نیوتن ضروری است.

یکی از مهم‌ترین محدودیت‌های روش نیوتن، نیاز به محاسبه مشتق تابع است. اگر تابع مشتق‌پذیر نباشد یا محاسبه مشتق آن دشوار باشد، استفاده از روش نیوتن با مشکل مواجه می‌شود. در چنین حالتی، روش‌هایی مانند روش سکانت می‌توانند جایگزین مناسبی باشند، زیرا در آن‌ها مشتق تابع به صورت تقریبی محاسبه می‌شود.

محدودیت دوم، صفر شدن مشتق در یکی از مراحل تکرار است. همان‌طور که فرمول \eqref{eq:ch2-newton} نشان می‌دهد، اگر $f'(x_n)=0$ باشد، این فرمول تعریف‌نشده خواهد بود. حتی اگر مشتق دقیقاً صفر نباشد اما مقدار بسیار کوچکی داشته باشد، کسر $\frac{f(x_n)}{f'(x_n)}$ ممکن است بسیار بزرگ شود و تقریب بعدی را از ریشه دور کند.

محدودیت سوم، وابستگی روش به حدس اولیه است. روش نیوتن یک روش محلی محسوب می‌شود؛ یعنی معمولاً زمانی عملکرد خوبی دارد که حدس اولیه به اندازه کافی به ریشه واقعی نزدیک باشد. اگر $x_0$ نامناسب انتخاب شود، روش ممکن است به ریشه مورد نظر همگرا نشود یا اصلاً همگرا نشود.

به عنوان مثال، اگر تابع دارای چند ریشه باشد، انتخاب حدس اولیه‌های مختلف می‌تواند منجر به همگرایی به ریشه‌های متفاوت شود. همچنین در برخی توابع، ممکن است دنباله تولیدشده توسط روش نیوتن بین چند مقدار نوسان کند و به هیچ ریشه‌ای نرسد.

محدودیت دیگر مربوط به ریشه‌های چندگانه است. اگر $\alpha$ ریشه‌ای چندگانه باشد، یعنی علاوه بر $f(\alpha)=0$، مشتق تابع نیز در آن نقطه صفر شود، سرعت همگرایی روش نیوتن کاهش می‌یابد. در چنین مواردی، همگرایی ممکن است از درجه دوم به درجه خطی کاهش پیدا کند.

به طور خلاصه، برخی از حالت‌های شکست یا ضعف روش نیوتن عبارت‌اند از:
\begin{itemize}
    \item صفر شدن مشتق در یکی از نقاط تکرار؛
    \item بسیار کوچک بودن مقدار مشتق و ایجاد گام‌های بزرگ؛
    \item انتخاب حدس اولیه نامناسب؛
    \item وجود چند ریشه و همگرایی به ریشه ناخواسته؛
    \item نوسان میان چند مقدار و عدم همگرایی؛
    \item کاهش سرعت همگرایی در ریشه‌های چندگانه؛
    \item مشتق‌ناپذیر بودن تابع در نزدیکی ریشه.
\end{itemize}

بنابراین، اگرچه روش نیوتن بسیار قدرتمند است، استفاده موفق از آن نیازمند توجه به ویژگی‌های تابع، انتخاب حدس اولیه مناسب و بررسی معیارهای توقف است.

\section{جمع‌بندی فصل}

در این فصل، مبانی نظری روش نیوتن برای حل معادلات غیرخطی بررسی شد. ابتدا مفهوم ریشه تابع و مسئله ریشه‌یابی معرفی شد. سپس توضیح داده شد که در بسیاری از مسائل عملی، یافتن جواب دقیق ممکن نیست و به همین دلیل از تقریب‌های عددی استفاده می‌شود.

در ادامه، مفهوم خطا و معیارهای توقف در روش‌های عددی مورد بررسی قرار گرفت. معیارهایی مانند اختلاف دو تقریب متوالی (\eqref{eq:ch2-stop-diff2}) و مقدار تابع در نقطه تقریبی (\eqref{eq:ch2-stop-func2})، از جمله معیارهای پرکاربرد برای توقف الگوریتم هستند.

سپس ایده هندسی روش نیوتن توضیح داده شد. در این روش، در هر مرحله خط مماس بر نمودار تابع در نقطه فعلی رسم می‌شود و محل برخورد این خط با محور افقی به عنوان تقریب بعدی ریشه انتخاب می‌گردد. بر اساس این ایده، فرمول اصلی روش نیوتن (\eqref{eq:ch2-newton}) به دست آمد.

همچنین نشان داده شد که این فرمول تنها زمانی قابل استفاده است که $f'(x_n)\neq 0$ باشد. سپس روش نیوتن از دیدگاه بسط تیلور نیز تفسیر شد و مشخص گردید که این روش بر پایه تقریب خطی تابع بنا شده است.

در بخش همگرایی، بیان شد که اگر تابع شرایط مناسبی داشته باشد و حدس اولیه به ریشه نزدیک باشد، روش نیوتن دارای همگرایی درجه دوم است (\eqref{eq:ch2-quadratic-error}). این ویژگی باعث می‌شود روش نیوتن نسبت به بسیاری از روش‌های عددی دیگر سریع‌تر به جواب برسد.

در پایان، محدودیت‌ها و حالت‌های شکست روش نیوتن بررسی شد. مواردی مانند صفر شدن مشتق، انتخاب حدس اولیه نامناسب، وجود ریشه‌های چندگانه و مشتق‌ناپذیری تابع می‌توانند عملکرد روش را تحت تأثیر قرار دهند.

مطالب این فصل پایه نظری لازم برای فصل بعد را فراهم می‌کند. در فصل بعد، الگوریتم روش نیوتن به صورت مرحله‌به‌مرحله، همراه با شبه‌کد و فلوچارت ارائه خواهد شد.